\documentclass[11pt,a4paper]{article}
\usepackage[utf8]{inputenc}
\usepackage[T1]{fontenc}
\usepackage[portuguese]{babel}
\usepackage{xcolor}
\usepackage{hyperref}
\usepackage{geometry}
\usepackage{tcolorbox}

\geometry{margin=2.5cm}

\definecolor{primary}{RGB}{41, 128, 185}
\definecolor{secondary}{RGB}{44, 62, 80}
\definecolor{codebg}{RGB}{240, 240, 240}

\hypersetup{colorlinks=true, linkcolor=primary, urlcolor=primary}

\title{\color{secondary}\Huge\textbf{Exercício 5.2: Primeiros passos com a malha de serviço Istio}}
\author{\color{primary}DevOps com Kubernetes -- 2026}
\date{}

\begin{document}
\maketitle

\section*{\color{primary}Instruções do Exercício}
\begin{tcolorbox}[colback=blue!5!white,colframe=blue!75!black,title=Exercício 5.2: Primeiros passos com a malha de serviço Istio]
Aprofunde-se no ecossistema Service Mesh.

- Leia a documentação oficial sobre os fundamentos do Istio.

- Instale a ferramenta Istio CLI na sua máquina e configure o Istio no seu cluster k3d, lembrando-se das pré-condições.

- Faça o \textit{deploy} da aplicação de exemplo oficial (Sample app) no \textit{namespace} para visualizar o fluxo de tráfego, e siga os passos do tutorial até a etapa de 'Clean up' (limpeza).

- Certifique-se de que o Prometheus está instalado no cluster para coletar os dados do Kiali e do Istio.
\end{tcolorbox}

\end{document}
